\documentclass[11pt]{article}
\usepackage[preprint]{acl}
\usepackage{times}
\usepackage{latexsym}
\usepackage[T1]{fontenc}
\usepackage[utf8]{inputenc}
\usepackage{microtype}
\usepackage{inconsolata}
\usepackage{graphicx}
\usepackage{booktabs}
\usepackage{amsmath}
\usepackage{amssymb}
\usepackage{array}
\usepackage{tabularx}
\usepackage{enumitem}
\usepackage{xspace}
\setlength\titlebox{7.0cm}

\newcommand{\modelmain}{XLM-R large\xspace}
\newcommand{\modelaux}{mBERT\xspace}

\title{Cross-Lingual Concept Transfer via Representation Engineering\\Mid-Project Report}

\author{Sudarshan Nikhil \\
  IIIT Hyderabad \\
  \texttt{sudarshan.nikhil@students.iiit.ac.in}
  \And
  Agyeya Negi \\
  IIIT Hyderabad \\
  \texttt{agyeya.negi@students.iiit.ac.in}
  \And
  Harshikaa Agrawal \\
  IIIT Hyderabad \\
  \texttt{harshikaa.agrawal@students.iiit.ac.in}
  \And
  Amey Bangera \\
  IIIT Hyderabad \\
  \texttt{amey.bangera@students.iiit.ac.in}}

\begin{document}
\maketitle

\begin{abstract}
We study whether an abstract concept direction learned from English can be used to \emph{causally steer} a multilingual encoder in other languages. Our target setting is deliberately narrower and more mechanistic than standard multilingual evaluation: instead of asking only whether a model transfers, we ask whether internal concept structure is aligned enough to be \emph{intervention-ready}. Concretely, we focus on moral and social concepts, use English as the source language, Hindi as a high-coverage non-English language, and a configurable lower-resource language (Welsh by default in the present codebase). The mid-project milestone has prioritized building a complete, reproducible pipeline before launching the expensive full experimental sweep. We therefore report the project motivation, updated related work, finalized data and model choices, a complete methodology, the evaluation protocol, and detailed implementation progress. Our current contribution at the mid stage is an end-to-end research pipeline covering dataset construction, representation caching, alignment diagnostics, probe transfer, concept vector extraction, causal intervention hooks, ablation controls, and paper-ready plotting. The remaining work is to execute the full runs, verify a gold subset manually, and analyze the resulting failure modes for the final report and viva.
\end{abstract}

\section{Introduction}
Multilingual encoders often exhibit impressive zero-shot transfer for concrete lexical and syntactic phenomena, but their behavior becomes less stable when the target property is abstract, culturally loaded, or weakly grounded in surface form. Our project asks whether such instability is merely a performance symptom or whether it reflects a deeper problem in the geometry of multilingual representations. If abstract concepts are only partially aligned across languages, then the model may ``know'' the concept in English while failing to access it reliably in Hindi or another lower-resource language. Conversely, if cross-lingual abstract structure is sufficiently shared, then an English-derived concept direction should be usable as an intervention for non-English reasoning.

This distinction matters for both science and practice. From a scientific perspective, it lets us move from descriptive multilingual evaluation to a causal question about the usability of internal representations. From a practical perspective, it matters for safety and fairness: if multilingual models encode different or unstable moral tendencies across languages, then downstream systems may behave inconsistently in culturally important settings \citep{hammerl2023moralbias,agarwal2024ethical,khandelwal2024dit}. Our proposal therefore combines cross-lingual alignment diagnostics with representation engineering \citep{zou2023repe,turner2023actadd}.

The mid-semester requirement is an 8--10 page report that clarifies the project direction and the progress so far. The main message of this report is simple: between the initial proposal and the mid submission, we prioritized finishing the research pipeline end-to-end instead of presenting early but fragile numbers. This choice is deliberate. Because the eventual experiments involve multiple models, languages, layer sweeps, intervention-strength sweeps, and ablation controls, a partial pipeline would have produced results that were difficult to trust, reproduce, or defend in a viva. Our goal at the mid stage is therefore to ensure that once the compute jobs are started, they immediately yield analysable, publication-ready outputs.

\section{Problem Definition and Mid-Project Scope}
Let a multilingual encoder $f_{\theta}$ map an input sequence $x$ to layer-wise hidden states $\{h^{(\ell)}(x)\}_{\ell=1}^{L}$. After pooling, each layer gives a representation
\begin{equation}
    r^{(\ell)}(x) = \mathrm{Pool}(h^{(\ell)}(x)) \in \mathbb{R}^{d}.
\end{equation}
For a concept $c$ such as \textsc{care/harm} or \textsc{fairness/cheating}, we estimate an English concept direction $v_c^{(\ell)}$ at layer $\ell$. Our main causal test is whether adding this direction inside the forward pass of a non-English example changes a downstream moral prediction in the expected direction:
\begin{equation}
    \tilde{h}^{(\ell)}(x) = h^{(\ell)}(x) + \alpha v_c^{(\ell)}.
\end{equation}

The project is organized around three questions:
\begin{enumerate}[leftmargin=*,itemsep=2pt]
    \item \textbf{RQ1 (alignment):} How aligned are English, Hindi, and lower-resource representations for abstract moral concepts?
    \item \textbf{RQ2 (controllability):} Does an English-derived concept direction reliably steer non-English predictions?
    \item \textbf{RQ3 (failure modes):} When steering fails, is the dominant cause layer choice, language resource level, translation/tokenization quality, or deeper semantic mismatch?
\end{enumerate}

At the mid-project stage, our scope is to finalize the data pipeline, models, alignment metrics, intervention machinery, and reporting infrastructure. Full result tables and qualitative case studies are scheduled for the next milestone after the large runs complete.

\begin{figure*}[t]
\centering
\fbox{\parbox{0.97\textwidth}{
\textbf{Implemented pipeline overview.} (1) Build a balanced English seed set from \textsc{MoralTextManipulation}; translate it into Hindi and a configurable lower-resource language; export a manually checkable gold subset. (2) Cache layer-wise encoder representations for \modelmain and \modelaux in HDF5. (3) Measure cross-lingual alignment with cosine similarity, linear CKA, and English-to-non-English probe transfer. (4) Extract concept directions per layer using mean-difference, one-vs-rest probe weights, and PCA-residual methods. (5) Run causal steering by injecting an English concept direction inside the encoder forward pass and reading out the effect with a trained linear probe. (6) Run random-direction and shuffled-concept ablations, then correlate alignment with intervention effect and generate publication-ready plots.
}}
\caption{The mid-project codebase already covers the full experimental pipeline. The remaining work is compute execution, manual verification of a gold subset, and final analysis of the resulting outputs.}
\label{fig:pipeline}
\end{figure*}

\section{Background and Related Work}
\paragraph{Multilingual encoders and transfer.}
Multilingual BERT and XLM-R are standard reference points for cross-lingual transfer \citep{devlin2019bert,conneau2020xlmr}. XLM-R is particularly relevant here because it improves multilingual performance at scale and is notably stronger than earlier multilingual baselines on several transfer tasks \citep{conneau2020xlmr}. Our codebase therefore uses \modelmain as the main model and \modelaux as a secondary baseline for model-size and architecture comparison.

\paragraph{Cross-lingual moral and ethical variation.}
Several recent papers show that moral or ethical judgments in language models can vary substantially with prompt language and may not track human cultural distributions cleanly. \citet{hammerl2023moralbias} adapt the \textsc{MORALDIRECTION} framework to multilingual settings and show that multilingual models encode different moral biases across languages. \citet{agarwal2024ethical} and \citet{khandelwal2024dit} both report language-sensitive changes in ethical reasoning and moral judgment, including systematic drops for Hindi and some lower-resource languages. These works motivate our problem setting but remain mostly diagnostic: they expose inconsistency without testing whether the internal representation can be causally controlled across languages.

\paragraph{Representation engineering and activation steering.}
Representation engineering treats population-level directions in hidden-state space as meaningful units for analysis and control \citep{zou2023repe}. Activation engineering methods such as ActAdd show that inference-time modification of internal activations can steer high-level behavior without model retraining \citep{turner2023actadd}. These ideas motivate our causal intervention stage. However, almost all prominent demonstrations of activation steering are English-centric and mostly focused on decoder-style models, while our project studies multilingual encoders and concept transfer across languages.

\paragraph{Controlling information in representations.}
Work such as INLP demonstrates that linear structure in hidden states can be manipulated for downstream control and analysis \citep{ravfogel2020null}. Although the goal there is information removal rather than positive steering, the underlying assumption is shared with our project: if a concept is linearly recoverable in representation space, interventions based on that structure may reveal how robust or language-specific the encoding is.

\paragraph{Updated data and benchmarking context.}
The specific dataset used in our implementation, \textsc{MoralTextManipulation}, was published in 2025 and provides large-scale moral annotations over both human-authored and LLM-manipulated texts \citep{greco2025moraltext}. Newer work such as \textsc{MEXA} formalizes cross-lingual alignment as a predictor of multilingual model behavior \citep{kargaran2025mexa}, while \textsc{MFTCXplain} highlights that multilingual moral reasoning remains challenging, especially in underrepresented languages \citep{trager2025mftcxplain}. These works strengthen the case for treating alignment and controllability jointly rather than as separate phenomena.

\section{Data, Languages, and Annotation Strategy}
\subsection{Seed Dataset}
We use the \textsc{MoralTextManipulation} resource \citep{greco2025moraltext} as the English seed set. The implementation currently uses the \texttt{unconditioned}/\texttt{revise} portion of the dataset because it provides short, relatively clean English texts with moral labels that can be converted into a single-label six-way formulation: \textsc{CH} (care/harm), \textsc{FC} (fairness/cheating), \textsc{LB} (loyalty/betrayal), \textsc{AS} (authority/subversion), \textsc{PD} (purity/degradation), and \textsc{NM} (non-moral). We filter the source data to examples with a clean single-label interpretation and then draw a class-balanced sample.

\subsection{Language Design}
English is the source language for concept extraction. Hindi is the primary non-English target because it is directly relevant to our project motivation and to recent findings on multilingual moral instability \citep{agarwal2024ethical,khandelwal2024dit}. For the lower-resource condition, the current repository defaults to Welsh (\texttt{cy}) in order to keep the initial run reproducible with readily available OPUS-MT translation checkpoints \citep{tiedemann2020opusmt}. Importantly, the codebase is parameterized rather than hard-coded: the lower-resource language can be swapped without changing the rest of the pipeline.

\subsection{Translation and Quality Control}
The present implementation constructs a parallel-style dataset automatically using OPUS-MT translation models \citep{tiedemann2020opusmt}. This is efficient, but it introduces translation noise precisely where our concepts are subtle and culturally sensitive. To address this, our design separates \textbf{silver} data (machine-translated examples used for scaling the main experiments) from a \textbf{gold} subset that will be manually verified. The code already includes a utility for exporting this gold subset for human checking. In the final analysis, we will report all main numbers on the silver set and verify conclusions on the gold subset.


\subsection{Dataset Schema and Annotation Artifacts}
The processed dataset is stored in JSONL format so that every stage of the pipeline can consume the same canonical representation. Each record contains an example identifier, the English source text, translated Hindi and lower-resource variants, the collapsed six-way label, optional source metadata from the original dataset, and split information. We also keep language-specific text fields separate rather than storing all languages in a single nested blob because this makes representation caching, debugging, and error analysis easier.

This schema is intentionally designed for the final paper as much as for the experiments themselves. Once the runs complete, the same records can be joined with alignment scores, probe predictions, intervention outputs, and qualitative notes, giving us a unified artifact for later analysis. In other words, we are not merely preparing inputs; we are building the eventual analysis table in a form that can support error slices by concept, language, layer, or intervention method.

\begin{table}[t]
\centering
\small
\begin{tabular}{p{0.18\columnwidth}p{0.22\columnwidth}p{0.20\columnwidth}p{0.30\columnwidth}}
\toprule
\textbf{Item} & \textbf{Choice} & \textbf{Role} & \textbf{Why this choice} \\
\midrule
Source dataset & MoralTextManipulation & English seed data & Large, recent, moral labels already available \\
Main model & XLM-R large & Primary encoder & Strong multilingual transfer baseline \\
Secondary model & mBERT & Comparison encoder & Widely used multilingual baseline \\
Target language 1 & Hindi & Main non-English language & High relevance and known multilingual moral variation \\
Target language 2 & Welsh (default) & Lower-resource condition & Runnable now; code remains language-configurable \\
Translation & OPUS-MT & EN$\rightarrow$HI/LR generation & Open checkpoints and easy integration \\
Data split & 80/10/10 & Train/dev/test & Stable probe fitting and evaluation \\
\bottomrule
\end{tabular}
\caption{Finalized mid-project choices for the first full experimental run.}
\label{tab:choices}
\end{table}

\section{Methodology}
\subsection{Representation Caching}
For each text in each language, we run the encoder once and cache pooled hidden states for every layer. This choice is central to the project design: once the HDF5 cache exists, all alignment diagnostics, probe fitting, concept extraction, and many analysis passes become much cheaper and more reproducible. The cache stores arrays with shape $(N, L, D)$ for each split and language.

We support both \textsc{CLS} pooling and masked mean pooling, although the current default is \textsc{CLS} for consistency across modules. The cache format also records metadata such as model ID, number of layers, representation dimensionality, pooling choice, and random seed.

\subsection{Alignment Diagnostics}
We measure alignment in three complementary ways.
\begin{enumerate}[leftmargin=*,itemsep=2pt]
    \item \textbf{Pairwise cosine alignment.} For parallel or comparable instances, we measure the cosine similarity between English and non-English representations at each layer.
    \item \textbf{Linear CKA.} Because simple pairwise cosine can miss structure at the matrix level, we also compute linear CKA \citep{kornblith2019cka} between representation matrices for each language pair and layer.
    \item \textbf{Probe transfer.} We train a multinomial probe on English train representations and test it on English, Hindi, and the lower-resource language layer by layer.
\end{enumerate}
Together, these metrics let us ask whether a layer is geometrically aligned, task-useful, both, or neither.

\subsection{Concept Direction Extraction}
For each label $c$ and layer $\ell$, we estimate an English concept direction using three methods.
\paragraph{Difference of means.}
We take the normalized difference between the mean English positive representation and the mean negative representation:
\begin{equation}
    v_{c,\mathrm{mean}}^{(\ell)} = \frac{\mu_{c,+}^{(\ell)} - \mu_{c,-}^{(\ell)}}{\lVert \mu_{c,+}^{(\ell)} - \mu_{c,-}^{(\ell)} \rVert_2}.
\end{equation}
\paragraph{Probe-weight direction.}
We fit a one-vs-rest linear probe for each label and use the normalized probe weight vector as the concept direction.
\paragraph{PCA-residual direction.}
We center the positive and negative groups, stack the resulting residuals, fit one-component PCA, and orient the principal component so that it points from negatives toward positives.

Using all three methods is important because a single extraction scheme could accidentally overfit to dataset artifacts. Agreement across methods would strengthen our claims; disagreement would itself be informative.

\subsection{Causal Steering in an Encoder-Only Setting}
A key implementation decision was required here. Our initial code path used a generic sequence-classification wrapper with an automatically resized output layer. That design was easy to run but methodologically weak because the classification head was not trained for our six-way task. We therefore replaced it with a more faithful setup: the model remains a frozen encoder, and the downstream readout is a \emph{trained linear probe} fitted on cached English representations. Steering is applied through PyTorch forward hooks at a chosen layer, and the effect is read out through that fixed probe.

This design is much closer to the proposal itself. It keeps the intervention inside the encoder, avoids retraining the multilingual model, and gives a meaningful downstream quantity for $\Delta$logit, $\Delta p$, average treatment effect (ATE), and success rate.

\subsection{Controls, Ablations, and Analysis}
To avoid overclaiming from any observed steering effect, our codebase includes the following controls:
\begin{itemize}[leftmargin=*,itemsep=2pt]
    \item a \textbf{random-direction control}, which preserves magnitude but destroys concept content;
    \item a \textbf{shuffled-concept control}, which injects the wrong concept direction for the same target label;
    \item a \textbf{layer sweep}, which identifies where controllability is strongest;
    \item an \textbf{alpha sweep}, which checks monotonicity and saturation of the intervention effect; and
    \item an \textbf{alignment--controllability correlation} stage, which tests whether layers that are better aligned are also easier to steer.
\end{itemize}
These are essential because a nonzero intervention effect alone does not imply concept-specific control.


\subsection{Experimental Matrix and Planned Comparisons}
Our final experiment set is a Cartesian product of several axes: encoder (\modelmain versus \modelaux), target language (Hindi versus lower-resource), layer, concept-direction extraction method, and intervention strength. This may sound straightforward on paper, but in practice it is the main reason we prioritized engineering completeness before reporting numbers. If any one stage were unstable, the entire matrix would become difficult to interpret.

\begin{table}[t]
\centering
\small
\begin{tabular}{p{0.21\columnwidth}p{0.53\columnwidth}}
\toprule
\textbf{Axis} & \textbf{Values in current plan} \\
\midrule
Encoder & XLM-R large, mBERT \\
Languages & English, Hindi, lower-resource (Welsh by default) \\
Layers & Full layer sweep for diagnostics; focused sweep for interventions \\
Direction methods & Mean-difference, probe-weight, PCA-residual \\
Intervention strength & $\alpha \in \{0, 0.25, 0.5, 1.0, 2.0\}$ \\
Controls & Random direction, shuffled concept, within-language upper bound \\
\bottomrule
\end{tabular}
\caption{Experimental axes already supported by the current codebase.}
\label{tab:matrix}
\end{table}

This matrix directly supports several comparisons that we plan to foreground in the final report: whether stronger multilingual encoders are more intervention-ready, whether better aligned layers are also the best steering layers, whether concept direction quality matters more than raw alignment, and whether lower-resource transfer failure is caused by weak representations or by unstable readout geometry.

\section{Implementation Progress So Far}
Table~\ref{tab:progress} summarizes the actual mid-project status. The important point is that the project is not at the ideation stage anymore; it is at the pipeline-complete stage.

\begin{table*}[t]
\centering
\small
\begin{tabular}{p{0.19\textwidth}p{0.11\textwidth}p{0.26\textwidth}p{0.34\textwidth}}
\toprule
\textbf{Module} & \textbf{Status} & \textbf{Current outputs} & \textbf{Notes} \\
\midrule
Balanced parallel dataset builder & Completed & JSONL dataset, stratified splits, summary JSON & Filters clean single-label rows and translates EN into HI and LR \\
Gold-subset export for manual verification & Completed & CSV export utility & Will be used for human verification of a smaller evaluation subset \\
Representation caching & Completed & HDF5 caches for both encoders & Stores pooled hidden states for all layers and all splits \\
Alignment diagnostics & Completed & Cosine and CKA JSON outputs & Ready for full sweeps once caches are built \\
Cross-lingual probe transfer & Completed & Accuracy, macro-F1, per-class F1, confusion matrices & One English-trained probe per layer; evaluates EN/HI/LR transfer \\
Concept vector extraction & Completed & Per-label, per-layer directions for 3 methods & Mean-diff, probe-weight, and PCA-residual variants \\
Encoder intervention hooks & Completed & Layer and alpha sweep scripts & Hook-based steering inside the encoder forward pass \\
Probe-based causal readout & Completed & Cached readout JSON + steering metrics & Replaced an earlier weaker design using an untrained classification head \\
Ablation controls & Completed & Random-direction and shuffled-concept JSON outputs & Designed to test concept specificity \\
Correlation analysis & Completed & Alignment-vs-ATE JSON and scatter plots & Tests whether aligned layers are also steerable \\
Figure generation & Completed & Mid-sem and post-midsem plotting scripts & Produces paper-ready plots from JSON outputs \\
Unit tests and run scripts & Completed & 25 passing tests and end-to-end shell scripts & Gives a stable base before expensive experiment execution \\
Full result sweeps on real data & Ongoing next milestone & Pending compute outputs & Scheduled immediately after report submission \\
Gold subset verification + qualitative analysis & Planned & Pending annotation and case studies & Needed for final interpretation and ethics discussion \\
\bottomrule
\end{tabular}
\caption{Mid-project progress summary. The heavy-lift experimental execution is the main remaining step, but the pipeline needed to produce trustworthy results is already in place.}
\label{tab:progress}
\end{table*}

In practical terms, the pipeline now runs from a single command and produces exactly the artifact types we want for the final evaluation: alignment curves, probe-transfer heatmaps, concept-vector files, intervention sweep JSONs, ablation outputs, and correlation plots. We also added tests and fixed a concrete evaluation bug in the classification helper so that metrics remain stable even when a split contains only a subset of labels.

It is equally important to state what is \emph{not} being claimed at the mid stage. We are not yet claiming any cross-lingual steering gain, any causal effect size, or any robust ranking between languages. Those statements depend on the full run. What we can claim now is narrower but still substantial: the project has crossed the difficult transition from idea to executable system. For a project like ours, that transition is the highest-risk part of the semester because failures in data design or intervention logic can invalidate later results. Reaching a stable pipeline before the compute-heavy phase therefore represents genuine scientific progress, not merely engineering housekeeping.

\section{Software Architecture and Reproducibility}
The repository is organized as a lightweight research framework rather than a collection of disconnected notebooks. Data construction, representation caching, evaluation, intervention, and plotting are separated into modules with command-line entry points. This matters because reproducibility failures in multilingual experiments often come from hidden preprocessing differences between scripts. By making every major stage consume standardized JSONL or HDF5 artifacts, we reduce that risk substantially.

We also made two engineering decisions specifically for the viva and final write-up. First, each major script emits structured JSON summaries in addition to plots. This allows us to regenerate figures or aggregate results without rerunning the underlying model. Second, the shell scripts at the repository root reflect the actual execution order of the full experiment, so that an evaluator can trace the pipeline from raw data construction all the way to final analysis outputs.

\begin{table*}[t]
\centering
\small
\begin{tabular}{p{0.18\textwidth}p{0.22\textwidth}p{0.46\textwidth}}
\toprule
\textbf{Layer of the system} & \textbf{Representative artifacts} & \textbf{Purpose in the research pipeline} \\
\midrule
Data layer & Processed JSONL, split files, translation exports, gold-subset CSV & Creates a stable multilingual dataset and preserves traceability back to the original English seed examples \\
Representation layer & HDF5 caches, metadata JSON, pooling configuration & Decouples expensive encoder inference from later probing and intervention analysis \\
Evaluation layer & Probe metrics, confusion matrices, cosine/CKA summaries & Measures whether cross-lingual structure is aligned and task-usable before intervention \\
Intervention layer & Concept vectors, readout weights, layer/$\alpha$ sweep outputs & Performs causal tests of whether English directions change non-English behavior \\
Analysis layer & Correlation summaries, scatter plots, aggregated case sheets & Connects alignment, controllability, and failure modes into publication-ready evidence \\
\bottomrule
\end{tabular}
\caption{How the repository is organized as a reproducible research system rather than a one-off notebook workflow.}
\label{tab:arch}
\end{table*}

A useful way to view the codebase is as a five-layer architecture, summarized in Table~\ref{tab:arch}. The data layer prepares the multilingual dataset and exports human-verification slices. The representation layer then converts that dataset into dense per-layer caches that all later modules share. The evaluation layer answers the descriptive question ``what is aligned?'' before the intervention layer asks the causal question ``what can be controlled?'' Finally, the analysis layer turns all of these intermediate outputs into interpretable plots and error slices.

This layered design matters because it lets us isolate bugs and methodological changes. For example, the earlier issue with the intervention readout could be fixed entirely inside the intervention layer without forcing us to rebuild the dataset or representation caches. Similarly, if we later decide to swap Welsh for another lower-resource language, the rest of the stack remains almost unchanged apart from translation and metadata regeneration. That modularity is exactly what we want in a mid-project system that still needs to evolve over the next few weeks.

Reproducibility is enforced through fixed random seeds, saved configuration values, deterministic split generation, and explicit logging of model IDs, pooling modes, layers, and intervention strengths. While these choices may appear mundane compared with the conceptual contribution, they are precisely what turn a promising idea into a defensible experimental system.

\section{Planned Baselines, Comparisons, and Error Analysis}
A strong final paper will depend not only on whether steering ``works,'' but on whether we compare it against the right baselines. Our project therefore defines a hierarchy of comparisons before the full runs begin. The first baseline is the \textbf{no-intervention} setting, which tells us the ordinary transfer quality of an English-trained readout. The second is the \textbf{within-language steering upper bound}: if Hindi-derived vectors steer Hindi reliably but English-derived vectors do not, then the problem is likely cross-lingual transfer rather than representational weakness in the target language itself. The third and fourth baselines are the random-direction and shuffled-concept controls already discussed above, which test whether any observed effect is concept-specific rather than merely caused by perturbing the hidden state.

We also plan to compare representation-level diagnostics against intervention outcomes rather than treating them as separate stories. A useful layer for probe transfer may not be the best layer for steering, and that discrepancy would itself be a result. One plausible outcome is that middle layers are most language-aligned while later layers are most task-discriminative, creating a trade-off between transferability and steerability. Another plausible outcome is that one concept, such as fairness, transfers better than another, such as purity, because its decision boundary is less culturally entangled. These comparisons are central to the final narrative and are already supported by the current implementation.

\begin{table*}[t]
\centering
\small
\begin{tabular}{p{0.17\textwidth}p{0.24\textwidth}p{0.25\textwidth}p{0.24\textwidth}}
\toprule
\textbf{Failure mode} & \textbf{Likely symptom in metrics} & \textbf{How we will diagnose it} & \textbf{Why it matters scientifically} \\
\midrule
Translation distortion & Probe transfer drops only on translated languages; manual inspection reveals semantic drift & Compare silver and gold subsets; inspect translated examples by concept & Distinguishes multilingual representation failure from data-construction failure \\
Tokenization or script fragmentation & Strong degradation concentrated in one language and in specific layers & Inspect token lengths, language-specific errors, and layer-wise curves & Tests whether lower-resource issues are architectural rather than conceptual \\
Concept overlap & Confusions between related labels such as care/harm and fairness/cheating & Use confusion matrices, per-class F1, and intervention direction mismatch & Reveals whether the concept space is separable enough for steering \\
Over-steering or saturation & Improvement at small $\alpha$ but instability or collapse at large $\alpha$ & Analyze monotonicity curves and success-rate change across strengths & Shows whether direction quality is robust or brittle \\
Cultural non-isomorphism & Good translation quality but persistent disagreement under intervention & Examine gold subset and qualitative notes for culturally grounded cases & Prevents us from mislabeling principled disagreement as model error \\
\bottomrule
\end{tabular}
\caption{Failure modes that the final analysis is designed to separate.}
\label{tab:failures}
\end{table*}

Our qualitative analysis plan is similarly fixed in advance. For each major concept and language, we will save representative cases in four buckets: successful transfer, baseline failure but intervention success, intervention failure despite strong alignment, and clear cultural or translation ambiguity. This is important because a project on abstract concepts can easily become shallow if it reports only aggregate metrics. The qualitative component is where we can connect hidden-state geometry to concrete multilingual examples and explain \emph{why} a particular intervention helped or failed.

A final comparison we intend to highlight is model sensitivity. XLM-R large and mBERT differ not only in scale but also in pretraining data and expected multilingual robustness. If the stronger model is better aligned yet not much easier to steer, that would suggest that better multilingual representations do not automatically yield better intervention readiness. If both models show the same failure patterns, the result would indicate a more general property of multilingual encoders rather than a quirk of one architecture.

\section{Evaluation Plan and Why Results Are Deferred to the Next Milestone}
Our final quantitative evaluation will report three groups of metrics.
\begin{table}[t]
\centering
\small
\begin{tabular}{p{0.22\columnwidth}p{0.52\columnwidth}}
\toprule
\textbf{Family} & \textbf{Metrics} \\
\midrule
Alignment & Pairwise cosine by layer, linear CKA by layer, EN$\rightarrow$HI/LR probe transfer \\
Task-level readout & Accuracy, macro-F1, per-class F1, confusion matrices \\
Causal steering & Mean $\Delta$logit, mean $\Delta p$, ATE, success rate, monotonicity over $\alpha$ \\
\bottomrule
\end{tabular}
\caption{Planned evaluation outputs.}
\label{tab:metrics}
\end{table}

We do \textbf{not} include final numerical results in this report for a methodological reason rather than for lack of preparation. The computationally meaningful part of this project begins only after the full chain is stable: if data construction, probing, steering, ablations, and plotting are not already locked down, then any early numbers can change for implementation reasons instead of scientific ones. Our choice was therefore to freeze the pipeline first and only then start the large sweeps.

This is especially important because our project is closer to a mini research system than to a single notebook experiment. A full run includes two encoders, three languages, layer-wise diagnostics, multiple concept-extraction methods, at least one target concept for steering, alpha sweeps, and control experiments. The value of the current milestone is that this entire chain is now built and testable. Immediately after the mid submission, the same code can be launched to produce the result tables that we will discuss in the viva.


Our primary hypotheses for the next stage are therefore already fixed. We expect: (i) middle layers to be most aligned across languages, (ii) the best steering layers to overlap only partially with the best probe-transfer layers, (iii) within-language steering to form an upper bound over cross-lingual steering, and (iv) lower-resource transfer to degrade more sharply than Hindi when the concept is culturally loaded or translation-sensitive. Stating these hypotheses now is useful because it prevents retrofitting the narrative after results arrive.

We also plan to report negative outcomes explicitly. If English directions fail to transfer despite good probe transfer, that would suggest that recoverable task information and intervention-ready geometry are different things. If they transfer only under one extraction method, that would narrow the claim from ``abstract concept directions transfer'' to ``certain operationalizations of concept directions transfer.'' Either outcome would still be scientifically meaningful.

\section{Planned Outputs for the Viva and Final Report}
An important practical objective of the current codebase is that it should not merely produce raw metric dumps. It should produce \emph{analysis-ready} outputs that can be discussed clearly in a viva and then integrated into the final paper. For that reason, every major stage of the pipeline is tied to a specific output artifact. Alignment diagnostics produce per-layer JSON summaries and line plots. Probe-transfer evaluation produces confusion matrices and per-class scores. Intervention sweeps produce tables over layer and $\alpha$, along with aggregated effect sizes and success-rate summaries. The correlation stage then links these outputs into a final set of cross-panel figures.

\begin{table}[t]
\centering
\small
\begin{tabular}{p{0.28\columnwidth}p{0.44\columnwidth}}
\toprule
\textbf{Output} & \textbf{How it will be used} \\
\midrule
Layer-wise alignment plots & Show where multilingual geometry is shared or diverges \\
Probe-transfer heatmaps & Compare model/language usability without intervention \\
Intervention sweep tables & Demonstrate whether control is monotonic and concept-specific \\
Ablation comparison plots & Rule out random-perturbation explanations \\
Qualitative case sheets & Support discussion of translation and cultural mismatch \\
\bottomrule
\end{tabular}
\caption{Planned viva-ready outputs from the finished pipeline.}
\label{tab:vivaoutputs}
\end{table}

This output-oriented design is also why we considered the mid milestone successful even without final metrics. By the time of the viva, it is much more valuable to have a coherent set of figures whose provenance is clear than to have a few ad hoc accuracy numbers from a partially debugged system. In effect, the engineering work completed for this submission is what enables the next three days of computation to be immediately useful. The report you are reading therefore reflects not an absence of progress but a specific sequencing decision: build the measurement instrument first, then trust the readings it produces.

From a paper-writing perspective, this structure also reduces the risk of narrative drift. Because each output corresponds to one research question, the eventual final write-up can be organized cleanly: first establish whether alignment exists, then show whether aligned directions are usable for steering, then analyze where the picture breaks. The present report already mirrors that structure so that the transition from mid submission to final paper is incremental rather than a complete rewrite.

\section{Timeline from Mid Submission to Final Analysis}
The project plan after this report is now concrete.
\begin{itemize}[leftmargin=*,itemsep=2pt]
    \item \textbf{March 8--9, 2026:} freeze codebase, submit mid report, and launch the first full \modelmain run;
    \item \textbf{March 9--11:} generate alignment and probe-transfer outputs, then run the first intervention sweeps for Hindi;
    \item \textbf{March 11--13:} run ablations and produce the first plot set for the viva;
    \item \textbf{March 13--20:} verify the gold subset manually and perform initial error analysis;
    \item \textbf{March 20 onward:} extend to more concepts and possibly a second lower-resource setting, refine the paper narrative, and prepare the final presentation.
\end{itemize}

This timeline is feasible precisely because the mid-stage engineering work is now largely complete.

\section{Risks, Limitations, and Ethics}
The biggest scientific risk is \textbf{cultural non-isomorphism}. An English concept direction may fail in another language not because the model is ``bad'' but because the concept boundary itself is culturally or linguistically different. We therefore intend to treat disagreement carefully rather than equating it automatically with error.

A second risk is \textbf{translation leakage}. Since the current large-scale pipeline is built on machine translation, some apparent cross-lingual effects could partly reflect translation artifacts. This is why we explicitly separate silver-scale runs from a gold manually verified subset.

A third risk is \textbf{over-interpreting steering}. If a hook changes a probe score, this does not prove that the model has human-like moral understanding. It shows that the direction is behaviorally usable with respect to our readout. Our claims in the final paper will therefore be limited to representational alignment, intervention effect, and robustness under controls.

Finally, the topic itself is sensitive. Steering moral judgments across languages can easily become normative or English-centric if framed carelessly. We will therefore present the project as an analysis of model internals and cross-lingual stability, not as a method for imposing one moral system on another.

\section{Conclusion}
This mid-project report presents a completed experimental pipeline for studying cross-lingual concept transfer via representation engineering. The project remains centered on the original question: can an abstract English concept direction causally steer multilingual encoder behavior in other languages? Since the proposal, we have translated that question into a concrete, testable system covering data preparation, alignment diagnostics, probe transfer, concept extraction, causal intervention, controls, and plotting. The most important mid-stage decision was to prioritize a full, defensible pipeline over premature numbers. As a result, the project is now ready for the compute-heavy phase in which the actual alignment and steering results will be generated and analyzed.

\bibliography{custom}

\end{document}
